\DocumentMetadata{lang=en-US,pdfstandard=ua-2,tagging=on}
\documentclass{ximera}

\usepackage{unicode-math}
\usepackage[points]{ximeraExam}

% Fill in the basic exam information here.
\examtitle{Exam Title}
\examcourse{Course Number}
\examdate{Month Day, Year}
\examversion{Form A}

% Student mode is the default; keep this explicit in a starter template.
\noprintanswers

% Adjust these only if you want different printed defaults.
\answerlineWidth{2in}
\answerlineAlign{right}
\answerlineInKey{hide}

% Example page furniture.
\runningheader{\examcourse}{\examtitle}{\examversion}
\runningfooter{\examdate}{}{Page~\thepage}

\begin{document}

% Use a custom cover so the grade table can appear before the questions.
\begin{examcover}
  \examcoverheading{\examtitletext}

  \medskip
  \noindent
  \examcoursetext
  \hfill
  \examversiontext
  \par

  \medskip
  \noindent
  \examdatetext
  \par

  \bigskip
  \noindent
  Name: \rule{3in}{0.4pt}

  \bigskip
  \begin{center}
    \gradetable[h][noearned][4]
  \end{center}
\end{examcover}

\begin{questions}

  % With no explicit question value, Question 1 inherits 5 points
  % from its two parts.
  \question
  Answer the following.

  \begin{parts}
    \part[2] Compute \(2+3\).

    \begin{solution}[1in]
      \(2+3=5\).
    \end{solution}

    \answerline{Answer:}

    \part[3] Explain your reasoning.

    \begin{solution}[\stretch{1}]
      Add two to three to obtain five.
    \end{solution}
  \end{parts}

  \newpage

  % An explicit question value overrides any inherited total.
  \question[4]
  Choose the correct statement.

  \begin{multipleChoice}[layout=horizontal,max-per-row=3]
    \choice{First option}
    \choice[correct]{Second option}
    \choice{Third option}
  \end{multipleChoice}

  \begin{solution}[1in]
    The second option is correct.
  \end{solution}

  \bonusquestion[2]
  Optional bonus question.

  \begin{solution}[1in]
    Put the bonus solution here.
  \end{solution}

\end{questions}

\end{document}
